% Set page size and margins
\usepackage[
  a4paper,
  top=2cm,
  bottom=3cm,
  left=2cm,
  right=2cm,
  marginparwidth=1.75cm,
  footskip=2.05cm,
]{geometry}

% Useful packages
\usepackage[export]{adjustbox}
\usepackage{amsmath}
\usepackage{array}
\usepackage{caption}
\usepackage[strict]{changepage}
\usepackage{enumitem}
\usepackage{etoolbox}
\usepackage{float}
\usepackage{fullwidth}
\usepackage{graphicx, trimclip}
\usepackage[colorlinks=true, allcolors=blue]{hyperref}
\usepackage{hyperref}
\usepackage[noautomatic, nonewpage]{imakeidx}
\usepackage{multicol}
\usepackage[super]{nth}
\usepackage{outlines}
\usepackage{paracol}
\usepackage[section]{placeins}
\usepackage{setspace}
\usepackage{stfloats}
\usepackage{subcaption}
\usepackage[usetransparent=false]{svg}
\usepackage{tabularx}
\usepackage[subfigure]{tocloft}
\usepackage{tikz}
\usepackage{titlesec}
\usepackage{verbatim}
\usepackage{varwidth}
\usepackage{wrapfig}
\usepackage[most]{tcolorbox}
\usepackage{catchfile}
\usepackage{xstring}
\newtcolorbox{scaledfigure}[1][]{height fill, space to=\myspace,#1}
\hypersetup{
  colorlinks=true,
  linkcolor=goldenbrown,
  filecolor=magenta,
  urlcolor=cyan,
  pdftitle={Rome Total War\: The Boardgame Rule Book},
  pdfpagemode=UseNone,
}
% Set the default spacing between paragraphs. Remove indentation.
\usepackage[skip=6pt, indent=0pt]{parskip}
\setstretch{1}

% Get version from env
% \getenv{variable_name} just prints the value
% \getenv[\macro]{variable_name} stores the value in \macro for reusability
\newcommand{\getenv}[2][]{%
  \CatchFileEdef{\value}{"|echo \$#2"}{\endlinechar=-1}%
  \if\relax\detokenize{#1}\relax\value\else\let#1\value\fi}

% Add dots to the table of contents
\renewcommand{\cftsecleader}{\cftdotfill{\cftsecdotsep}}
\renewcommand\cftsecdotsep{\cftdot}
\renewcommand\cftsubsecdotsep{\cftdot}

\captionsetup[figure]{labelformat=empty}
\captionsetup[subfigure]{labelformat=empty, singlelinecheck=off, justification=centering}
\usetikzlibrary{shadows, shadows.blur, calc, backgrounds}

\setlength{\columnsep}{1cm}
\newtoggle{printable}
\newtoggle{noartbackground}
\newtoggle{githubbuild}

% Variables
\def\_assets{assets}

\def\images{\_assets/images}
\def\layout{\_assets/layout}
\def\fonts{\_assets/fonts}

\def\repourl{https://github.com/fmestevez/rtw-bg-unofficial-rulebook}
\def\bggthreadurl{https://boardgamegeek.com/thread/3426675/rulebook-rewrite-project}

% Colors
\definecolor{amber}{rgb}{1.0, 0.49, 0.0}
\definecolor{antiquewhite}{rgb}{0.98, 0.92, 0.84}
\definecolor{arylideyellow}{rgb}{0.96, 0.89, 0.58}
\definecolor{cadmiumgreen}{rgb}{0.0, 0.42, 0.24}
\definecolor{darkcandyapplered}{rgb}{0.64, 0.0, 0.0}
\definecolor{goldenbrown}{rgb}{0.6, 0.4, 0.08}

\newcommand{\svg}[2][10]{%
  {\raisebox{-0.15\height}{\includesvg[height=#1px]{\svgs/\detokenize{#2}.svg}}}%
}%
\newcommand{\svgunit}[2][10]{%
  {\raisebox{-0.1\height}{\includesvg[height=#1px]{\svgs/\detokenize{#2}.svg}}}%
}%
\newcommand{\svgeven}[2][10]{%
  \includesvg[height=#1px]{\svgs/\detokenize{#2}.svg}%
}%

% Command to frame images
\newcommand\framedimage[2][]{%
  \begin{tikzpicture}
    \draw (0, 0) node[inner sep=0] {\makebox[#1][c]{\includegraphics[width=#1]{#2}}};
    \draw [bordermidyellow, thick] ([xshift=+1pt, yshift=-1pt] current bounding box.north west) rectangle ([xshift=-1pt, yshift=1pt] current bounding box.south east);
    \draw [borderoutyellow, thick] (current bounding box.north west) rectangle (current bounding box.south east);
    \draw [borderinyellow, thick] ([xshift=+3pt, yshift=-3pt] current bounding box.north west) rectangle ([xshift=-3pt, yshift=3pt] current bounding box.south east);
  \end{tikzpicture}}
% End of drop frame definition

\titleformat{\section}
{\huge}
{\filright
\footnotesize
\enspace SECTION \thesection\enspace}
{8pt}
{\Huge\bfseries\filcenter\uppercase}

\newcommand{\addsection}[3][1]{
  \vspace*{-5em} % Adjust vertical positioning
  \hspace*{-1em}
  \makebox[0pt][l]{
    \raisebox{-\totalheight}[0pt][0pt]{
      \begin{tikzpicture}
        % Insert the banner image
        \node[anchor=north] (banner) at (0, 0) {\includegraphics[width=\columnwidth, keepaspectratio]{\layout/section_heading.png}};
        % Place the section title inside the banner
        \node[align=center, text=white, font=\fontsize{16}{16}\selectfont\bfseries] at ([yshift=-1cm] banner.center) {#2};
      \end{tikzpicture}
    }
  }
  % Reserve space below the banner to prevent overlap with body text
  \vspace{12em} % Adjust spacing to push body text further down
  \begin{fullwidth}[outermargin=0.16\textwidth, innermargin=0.16\textwidth]
    \begin{center}
      \vspace*{\lang_header_adjustment}
      \cleardoublepage\phantomsection\addcontentsline{toc}{section}{\protect\numberline{}\mbox{#2}}
    \end{center}
  \end{fullwidth}
  \vspace{1.75em}
}
%End of create section heading.

% Apply language-specific subsection spacings if defined
\ifdefined\subsectionspacing
  \subsectionspacing{}
\fi

\newcommand\picdims[4][]{%
  \setbox0=\hbox{\includegraphics[#1]{#4}}%
  \clipbox{.5\dimexpr\wd0-#2\relax{} %
    .5\dimexpr\ht0-#3\relax{} %
    .5\dimexpr\wd0-#2\relax{} %
    .5\dimexpr\ht0-#3\relax}{\includegraphics[#1]{#4}}}

\tikzset{
  thick/.style=      {line width=1.3pt},
  very thick/.style= {line width=1.7pt},
  ultra thick/.style={line width=2.2pt}
}

\definecolor{borderoutyellow}{HTML}{DBCA86}
\definecolor{borderinyellow}{HTML}{B09E69}
\definecolor{bordermidyellow}{HTML}{6f6749}
% Create note box
\providecommand{\notefont}[0]{\fontfamily{ptm}\selectfont}
\newcommand{\note}[2]{
  \begin{tikzpicture}
    \textlang{english}{\draw (0, 0) node[inner sep=0] {\makebox[\linewidth][c]{\picdims[width=\linewidth]{\linewidth}{#1\baselineskip}{\layout/table-background.jpg}}};}
    \draw [borderoutyellow, very thick] (current bounding box.north west) rectangle (current bounding box.south east);
    \draw [borderinyellow, thick] ([xshift=+2.8pt, yshift=-2.8pt] current bounding box.north west) rectangle ([xshift=-2.8pt, yshift=2.8pt] current bounding box.south east);
    \node at (current bounding box.center) {
      \begin{varwidth}{0.85\linewidth}
      \notefont{
        \color{arylideyellow}
        \hypersetup{linkcolor=amber}
        #2
        \hypersetup{linkcolor=goldenbrown}
      }
      \end{varwidth}
    };
    \begin{pgfonlayer}{background}
      \begin{scope}[blend mode=multiply]
        \draw [shade, blur shadow={shadow opacity=15}] (current bounding box.north west) rectangle (current bounding box.south east);
      \end{scope}
    \end{pgfonlayer}
  \end{tikzpicture}
}

\definecolor{darkcellborder}{HTML}{634831}
\definecolor{darkcellbg}{HTML}{20160C}

\newcommand{\darkcell}[2][0.9]{
  \begin{tikzpicture}
    \filldraw[line width=1.0pt, fill=darkcellbg, fill opacity=0.5, draw=darkcellborder] (0, 0) rectangle (\linewidth, #1);
    \node[text width=\linewidth, align=center] at (current bounding box.center) {\textbf{#2}};
  \end{tikzpicture}
}

\definecolor{lightcellborder}{HTML}{77543e}
\definecolor{lightcellbg}{HTML}{20160C}

\newcommand{\lightcell}[2][0.9]{
  \begin{tikzpicture}
    \filldraw[line width=1.0pt, fill=lightcellbg, fill opacity=0.25, draw=lightcellborder] (0, 0) rectangle (\linewidth, #1);
    \node[text width=\linewidth, align=center] at (current bounding box.center) {\color{white}#2};
  \end{tikzpicture}
}

% Commands to be used for automation generating printable version
\newcommand{\pagetarget}[2]{\label{#1}\hypertarget{#1}{#2}}
\newcommand{\pagelink}[2]{\hyperlink{#1}{#2}\iftoggle{printable}{ \textmd{(\pageshorthand\,\pageref{#1})}}{}}

% Stubs for polyglossia package used for Hebrew but not other languages
\providecommand{\mainlanguagename}[0]{}
\providecommand{\textlang}[2]{#2}

% Command for overlay circled text
\definecolor{goblin}{HTML}{3b7c33}
\newcommand\encircle[1]{%
  \tikz[baseline=(X.base)]
  \node (X) [draw=white, shape=circle, inner sep=0, fill=goblin, text=white, blur shadow={shadow blur steps=5}] {\strut \textbf{#1}};%
}

\providetoggle{textright2left}
% Background
\AddToHook{shipout/background}{%
  \iftoggle{noartbackground}{}{
    \put (0in,-\paperheight){\includegraphics[width=\paperwidth,height=\paperheight]{\layout/content.jpg}}
  }
  % \iftoggle{printable}{
  %   \iftoggle{textright2left}{
  %     \ifodd\value{page}
  %       \put (0in,-\paperheight){\includegraphics[width=\paperwidth,height=0.05\paperheight]{\layout/bottom-even.png}}
  %     \else
  %       \put (0in,-\paperheight){\includegraphics[width=\paperwidth,height=0.05\paperheight]{\layout/bottom-odd.png}}
  %   \fi}{%
  %     \ifodd\value{page}
  %       \put (0in,-\paperheight){\includegraphics[width=\paperwidth,height=0.05\paperheight]{\layout/bottom-odd.png}}
  %     \else
  %       \put (0in,-\paperheight){\includegraphics[width=\paperwidth,height=0.05\paperheight]{\layout/bottom-even.png}}
  %   \fi}
  % }{\put (0in,-\paperheight){\includegraphics[width=\paperwidth,height=0.05\paperheight]{\layout/bottom.png}}}
}

\makeindex[columns=3, title=,]

\begin{document}

% !TeX spellcheck = en_US
\thispagestyle{empty}
\begin{tikzpicture}[remember picture, overlay, inner sep=10pt]
  \iftoggle{noartbackground}{}{
    \node(cover)[anchor=center] at (current page.center) {
      \includegraphics[height=\paperheight, keepaspectratio]{\layout/cover.jpg}
    };
  }
  \node(title)[minimum width = \paperwidth, anchor=center, yshift=\dimexpr-10em\relax] at (current page.north) {
    \includegraphics[width=0.6\paperwidth]{\layout/cover_title.png}
  };
  % \node(subtitle)[anchor=center, yshift=10em] at (current page.south) {
  %   \includegraphics[width=0.6\paperwidth]{\layout/cover_subtitle.png}
  % };
\end{tikzpicture}


\iftoggle{printable}{
  \newgeometry{
    twoside,
    top=2cm,
    bottom=3cm,
    left=2.5cm,
    right=1.5cm,
    marginparwidth=1.75cm,
    footskip=2.05cm,
  }
}{}


% Let main_xx.tex provide translated \author
\ifdefvoid{\translatedauthor}{
  \author{\textlang{english}{Facundo Estevez}}
}{
  \author{\translatedauthor}
}
\maketitle

\begin{center}
  \iftoggle{githubbuild}{
    \getenv[\githubsha]{GITHUB_SHA}
    \versionwarning{} \href{\repourl}{\textlang{english}{\StrLeft{\githubsha}{7}}}.
  }{
    \textlang{\mainlanguagename}{\versionlabel{} \textlang{english}{\input{.version}}}
  }

  \bigbreak

  \intro{}

  % \iftoggle{printable}{}{\bigbreak\heegusquote}
\end{center}

\iftoggle{printable}{
  \bigbreak

  \begin{multicols}{2}
  \centering
  \includegraphics[width=0.8\linewidth]{\qr/github.png}\\
  \qrgithub

  \columnbreak

  \includegraphics[width=0.8\linewidth]{\qr/bgg.png}\\
  \qrbgg
  \end{multicols}
}{}

% \textlang{english}{\begin{tikzpicture}[remember picture, overlay]
%   \node(cover)[anchor=center, yshift=12em] at (current page.south) {
%     \includegraphics[width=1.01\paperwidth, keepaspectratio]{\art/castle_bottom.png}
%     \thispagestyle{empty}
%   };
% \end{tikzpicture}}

\clearpage

\begin{multicols*}{2}
\tableofcontents
\vspace*{\fill}
\columnbreak
\vspace*{\fill}
% \includegraphics[width=\linewidth]{\art/black_dragon.jpg}
\end{multicols*}

\clearpage

% !TeX spellcheck = en_US

\begin{multicols}{2}
\addsection{INTRODUCTION}

\iftoggle{printable}{}{\bigbreak}

The region surrounding the Mediterranean Sea has experienced conflict throughout most of its recorded history. During the period we know as the Ancient World, mighty empires rose and fell upon one another. Every empire’s rise and every empire’s fall were marked by war. Total War.

Based on the world-renowned video game, \pagetarget{Rome Total War}{\textbf{Rome Total War: The Board Game}} allows up to four players to play as the leaders of Rome, Carthage, Greece, and Gallo-Germania. Each player will claim territories, recruit armies, fight enemies, raise taxes, spy, riot, and more.

\textbf{Notice}: In this rule book, game and component terms are Capitalized.
\textbf{Bold text} is used to draw attention to important rules.
\textit{Italicization} is used for gameplay examples.
\pagelink{Rome Total War}{Brown colored hyperlinks} will take you to other parts of the rule book.
\vfill
\columnbreak
% \note{10}{Exceptions and notes with \pagelink{Rome Total War}{amber colored hyperlinks} are explained in boxes like this one.}
\vfill
\end{multicols}

% !TeX spellcheck = en_US

\restoregeometry

\pagestyle{empty}
\iftoggle{noartbackground}{}{
  \begin{tikzpicture}[remember picture, overlay, inner sep=10pt]
    \node(cover)[anchor=center] at (current page.center) {
      \includegraphics[height=\paperheight, keepaspectratio]{\layout/back.png}
    };
  \end{tikzpicture}
}


\end{document}
